\documentclass[11pt]{article}
\usepackage[T1]{fontenc}
\usepackage[margin=1in]{geometry}
\usepackage{booktabs}
\usepackage{array}
\newcolumntype{L}[1]{>{\raggedright\arraybackslash}p{#1}}

\title{ECEC state demand inputs: method}
\author{Anna Rachwalski}
\date{18 September 2026}

\pdfinfo{ /Author (Anna Rachwalski) /Title (ECEC state demand inputs: method) }

\begin{document}
\maketitle

The simulator needs one small file per geography: six rows, one per care type,
one number each, the annual hours of child care children under five use. This
memo says where that number comes from.

Everything happens at the state level. County data enters once, in the
population step, and it enters as a sum. Nothing joins on county.

\section{Why this is not a survey lookup}

The simulator calibrates each state against the mix of child care its families
actually use. NSECE, the survey that measures that mix, does not record which
state a family is in. It publishes Census region and nothing finer.

So I rebuild the mix from each state's own administrative records, and I use the
survey only for the things no state records can see.

\section{Results}

The build measures two things. Everything else is a subdivision of them.

\begin{center}
\begin{tabular}{L{3.2cm}rrrr}
\toprule
& \textbf{Center-based} & \textbf{Not center-based} & \textbf{Total} & \textbf{Share center} \\
\midrule
Kentucky      &  81,326 &  80,604 & 161,930 & 50\% \\
New Jersey    & 173,225 & 136,721 & 309,946 & 56\% \\
New York      & 270,982 & 378,847 & 649,829 & 42\% \\
New York City & 132,065 & 169,911 & 301,976 & 44\% \\
Virginia      & 135,068 & 127,932 & 263,000 & 51\% \\
\bottomrule
\end{tabular}
\end{center}

Care that is not center-based is one part counted and one part left over:

\begin{center}
\begin{tabular}{L{3.2cm}rr}
\toprule
& \textbf{Licensed home care} & \textbf{Informal, unrecorded} \\
\midrule
Kentucky      &  4,756 &  75,848 \\
New Jersey    &  1,987 & 134,734 \\
New York      & 96,366 & 282,481 \\
New York City & 38,642 & 131,269 \\
Virginia      &  4,743 & 123,189 \\
\bottomrule
\end{tabular}
\end{center}

New York's licensed home care is twenty times Kentucky's or Virginia's. That is
real rather than an artefact: New York registers home care from three children
up, so care that leaves no trace elsewhere sits inside its register.

Children, by care type. These are counts of children, not records and not
weighted survey figures:

\begin{center}
\begin{tabular}{L{4.0cm}rrrrr}
\toprule
& \textbf{KY} & \textbf{NJ} & \textbf{NY} & \textbf{NYC} & \textbf{VA} \\
\midrule
Center-Based        & 37,550 & 116,402 & 168,235 &  86,337 &  87,192 \\
Unpaid Center-Based & 43,776 &  56,823 & 102,747 &  45,728 &  47,876 \\
Paid Home-Based     &  4,756 &   1,987 &  96,366 &  38,642 &   4,743 \\
Other Unpaid        & 45,830 &  81,411 & 170,686 &  79,318 &  74,436 \\
Other Paid          & 27,363 &  48,607 & 101,909 &  47,357 &  44,442 \\
Unpaid Home-Based   &  2,655 &   4,715 &   9,886 &   4,594 &   4,311 \\
\midrule
\textbf{Total}      & \textbf{161,930} & \textbf{309,946} & \textbf{649,829}
                    & \textbf{301,976} & \textbf{263,000} \\
\bottomrule
\end{tabular}
\end{center}

Annual hours, in millions. This is the output file:

\begin{center}
\begin{tabular}{L{4.0cm}rrrrr}
\toprule
& \textbf{KY} & \textbf{NJ} & \textbf{NY} & \textbf{NYC} & \textbf{VA} \\
\midrule
Center-Based        & 63.3 & 196.2 & 283.6 & 145.5 & 147.0 \\
Unpaid Center-Based & 62.9 &  81.6 & 147.6 &  65.7 &  68.8 \\
Paid Home-Based     &  9.1 &   3.8 & 185.0 &  74.2 &   9.1 \\
Other Unpaid        & 88.4 & 157.1 & 329.3 & 153.0 & 143.6 \\
Other Paid          & 42.5 &  75.5 & 158.4 &  73.6 &  69.1 \\
Unpaid Home-Based   &  5.6 &  10.0 &  21.0 &   9.7 &   9.1 \\
\bottomrule
\end{tabular}
\end{center}

New York City is both inside New York and a build of its own. The two overlap
deliberately, and they do not match: each is calibrated against its own
non-parental total, so the five boroughs come out differently depending on which
build produced them. Read one or the other, not both.

\section{The build}

\subsection*{1. Count the children}

Census population estimates, vintage 2019, children aged 0 to 4, summed from
counties to the state.

\begin{center}
\begin{tabular}{L{3.2cm}r}
\toprule
& \textbf{Children 0--4, 2019} \\
\midrule
Kentucky      &   272,610 \\
New Jersey    &   514,690 \\
New York      & 1,127,001 \\
New York City &   523,718 \\
Virginia      &   505,477 \\
\bottomrule
\end{tabular}
\end{center}

I use ages 0 to 4 because the simulator does. Its ACS processing filters
\texttt{age < 5} in seven places, and the demand file has to describe the same
children the simulator is calibrating against. A 0--5 universe was built and
discarded: it raises every total by 15 to 23 percent, it needs a second
population source, and Census publishes no single-year-of-age county estimate,
so New York City could only be estimated rather than counted.

New York City is 47 percent of the state's young children.

\subsection*{2. How many are in someone else's care}

I multiply by a rate from a household survey. NSCH 2019 asks whether the child
was in any non-parental care in the past week, and it records the state.

\begin{center}
\begin{tabular}{L{3.2cm}rrrr}
\toprule
& \textbf{Children 0--4} & \textbf{Rate} & \textbf{In non-parental care} & \textbf{Records} \\
\midrule
Kentucky      &   272,610 & 59.4\% & 161,930 & 147 \\
New Jersey    &   514,690 & 60.2\% & 309,946 & 101 \\
New York      & 1,127,001 & 57.7\% & 649,829 & 112 \\
New York City &   523,718 & 57.7\% & 301,976 & 112 \\
Virginia      &   505,477 & 52.0\% & 263,000 & 121 \\
\bottomrule
\end{tabular}
\end{center}

That product is the control total. Every later count is a share of it.

The total is always 2019, because both inputs are 2019. That is what lets me use
a provider register published in 2026 without ageing it backward: the register
decides the mix, never the level.

New York City takes the state's rate, because NSCH publishes no sub-state
geography. Everything that distinguishes the city has to come from its
registers, and it does.

This rate is the softest number in the build. About 120 families per state,
which is roughly a seven point standard error on a figure that multiplies the
entire non-parental total. Section 5 says why I did not pool more waves to fix
it.

\subsection*{3. Turn provider registers into children}

Every state licenses child care providers and publishes the register. The
register gives capacity; I need children. Three things have to happen, and each
state needs them done differently because the registers differ.

\textbf{Remove the school-age capacity.} Every register mixes it in.

\begin{itemize}
\setlength{\itemsep}{0.2em}
\item \textbf{Kentucky} does not say which providers are centers and which are
homes. I split on the state's own rule: a licensed home is capped at 12
children, so capacity of 12 or less is a home. That gives 1,609 centers and 335
homes. Providers serving school-age children only are dropped. Mixed age bands
scale by ratios measured inside the file itself, 0.487 to 0.827 depending on the
band.
\item \textbf{New Jersey} publishes an age range per center. I drop the 1,035
centers serving only ages 6 to 13, and scale mixed ranges by 0.820.
\item \textbf{New York} is the clean case. It publishes capacity by age, so
infant plus toddler plus preschool is the under-five figure exactly. Homes
report one lump, but school-age capacity is published separately and subtracts.
\item \textbf{Virginia} writes the age range as text. I parse it to months and
drop anything starting at age five or later. Centers reaching past six scale by
0.60.
\end{itemize}

\textbf{Bring the register back to 2019.} Two of them are 2026 exports.

\begin{itemize}
\setlength{\itemsep}{0.2em}
\item \textbf{New York} scales by provider type against the state's own
2015--2025 capacity series. Family homes have closed in large numbers: group
family day care scales to 0.86 of its current size statewide and 0.84 in the
city. Centers barely moved, so one blanket factor would distort the mix.
\item \textbf{Virginia} scales by one statewide factor, 0.90, against its June
2020 capacity matrix.
\item \textbf{Kentucky and New Jersey} need nothing. Kentucky's register is a
September 2018 snapshot, ten months before the base year and the best vintage
here. New Jersey's is frozen at May 2019.
\end{itemize}

\textbf{Turn capacity into children.} Capacity is the legal maximum. Real
enrollment is lower.

\begin{itemize}
\setlength{\itemsep}{0.2em}
\item \textbf{Kentucky measured its own, and uses no fill rate at all.} Its 2019
workforce study prints capacity and enrollment side by side: 42 children per
center, 9 per home. This is the best way to do it, and it is why 1,609 centers
yield only 67,578 children.
\item \textbf{New Jersey measured its own too}, 67 percent full, from its 2021
market rate survey.
\item \textbf{New York and Virginia have published none}, so they borrow the
national figures: centers 90 percent full, homes 60 percent. New York's homes
take 80 percent instead, the fill rate without the age correction, because its
register already removed school-age capacity and applying the correction again
would remove it twice.
\end{itemize}

\begin{center}
\begin{tabular}{L{2.6cm}L{3.2cm}rrr}
\toprule
& & \textbf{Providers} & \textbf{Under-5 capacity} & \textbf{Children} \\
\midrule
Kentucky      & Centers        &  1,609 & 107,369 &  67,578 \\
Kentucky      & Homes          &    335 &   2,115 &   3,015 \\
New Jersey    & Centers        &  3,128 & 251,159 & 169,106 \\
New York      & Centers, state &  2,073 & 153,053 & 138,191 \\
New York      & Centers, city  &  2,310 & 141,300 & 127,580 \\
New York      & Homes          & 11,918 & 117,401 &  93,381 \\
New York City & Centers        &  2,310 & 141,300 & 127,580 \\
New York City & Homes          &  4,580 &  46,057 &  36,634 \\
Virginia      & Centers        &  3,186 & 207,835 & 187,654 \\
Virginia      & Homes          &  1,935 &  12,633 &   7,604 \\
\bottomrule
\end{tabular}
\end{center}

New York needs two registers, because it has two regulators. The state licenses
no centers inside New York City; the city permits its own under Health Code
Article 47. The state's row for the city is zero, and that is correct rather
than a join failure. Neither register alone covers the state.

\subsection*{4. Add the programs that count their own children}

Head Start, Early Head Start and public pre-K publish actual enrollment, so no
estimate is needed.

\begin{center}
\begin{tabular}{L{3.2cm}rrrr}
\toprule
& \textbf{Pre-K} & \textbf{Head Start} & \textbf{Early Head Start} & \textbf{Total} \\
\midrule
Kentucky      & 28,465 & 12,307 &  3,004 &  43,776 \\
New Jersey    & 41,305 & 11,940 &  3,578 &  56,823 \\
New York      & 54,451 & 36,863 & 11,433 & 102,747 \\
New York City & 26,717 & 13,073 &  5,938 &  45,728 \\
Virginia      & 33,790 & 11,579 &  2,507 &  47,876 \\
\bottomrule
\end{tabular}
\end{center}

These children come out of the licensed center count, not on top of it. A Head
Start center usually holds a state licence, so its children already sit inside
step 3's capacity, and adding them would count them twice.

A published count is treated as a census and does not move in step 5.

This category is an upper bound in every state. I cannot measure how much Head
Start overlaps state pre-K, and the published figures do not net it out.

\subsection*{5. Scale the level to the survey}

Provider registers are good at shape and bad at level. They say which care types
are common in a state; they overstate how many children are in formal care.

So I scale each geography by one factor until formal care matches the share the
survey finds, a little over half of all non-parental care. The program counts
hold still. Everything else moves together, so the mix is untouched and only the
level changes.

\begin{center}
\begin{tabular}{L{3.2cm}rrrr}
\toprule
& \textbf{Scaled by} & \textbf{Formal} & \textbf{Informal} & \textbf{Informal share} \\
\midrule
Kentucky      & 1.58 &  86,082 &  75,848 & 47\% \\
New Jersey    & 1.04 & 175,212 & 134,734 & 44\% \\
New York      & 1.03 & 367,348 & 282,481 & 44\% \\
New York City & 1.05 & 170,707 & 131,269 & 44\% \\
Virginia      & 0.62 & 139,811 & 123,189 & 47\% \\
\bottomrule
\end{tabular}
\end{center}

This table is the best diagnostic in the build.

New Jersey, New York and the city barely move. Their raw register counts landed
within two points of the survey before anything was done to them. Three separate
state registers agreeing with an independent national survey is the strongest
evidence here that the conversion in step 3 is sound.

Virginia is cut by a third. Its register is the most inclusive of the four,
listing unlicensed and voluntarily registered family homes that other states
never record, and it borrows the national fill rate.

Kentucky is raised by half, and that is the informative one. Kentucky is the
only state that measured how full its own centers are: 57 percent, against the
90 percent that New York and Virginia borrow. The borrowed figure is too high by
more than half. Kentucky moving up and Virginia moving down is exactly what that
gap predicts, so the scaling is correcting a bias two states have measured
rather than absorbing an unexplained one.

\subsection*{6. Finish the rows}

\textbf{There is no price split.} An earlier version divided center care into
low-priced and high-priced. Nothing in the state data supports that division. A
state publishes one price, so every state landed wholly on one side of the
national line, and a care type with no children in it reads as unmeasured rather
than as absent. Center care is now one row, at 1,686 hours per child per year,
the national average across both price bands.

\textbf{Fill in the care nobody records.} Three of the six types appear in no
register anywhere: unpaid relatives, nannies, informal arrangements. Unpaid
relative care is not a business and files no return. It is whatever is left of
the control total after everything countable is subtracted, divided three ways
on the national pattern. That remainder runs 44 to 47 percent of non-parental
care here, against 47 percent nationally.

\textbf{Convert children to hours.} Each care type has a measured hours figure
from the national survey: the child-weighted mean weekly hours in the child's
primary arrangement, times 50 weeks. They run from about 1,440 hours a year for
program center care to about 2,120 for unpaid home care, so the hours column is
not proportional to the children column.

\section{What each number rests on}

\begin{center}
\begin{tabular}{L{4.4cm}L{4.4cm}L{4.4cm}}
\toprule
& \textbf{Where it comes from} & \textbf{How solid} \\
\midrule
Children 0--4 & Census, 2019 & Counted \\
\addlinespace
Share in non-parental care & Household survey, 2019, about 120 families per
  state & Soft. Pooling was tested and rejected \\
\addlinespace
Provider capacity & State licensing register & Counted, but a legal maximum.
  Virginia providers say they can serve 27 percent fewer children than they are
  licensed for \\
\addlinespace
How full providers are, KY and NJ & The state's own survey & Measured in state \\
\addlinespace
How full providers are, NY and VA & National survey & Borrowed. Two states' own
  data say it is too high \\
\addlinespace
Program enrollment & Federal and state reporting & Counted \\
\addlinespace
Informal split, hours per child & National survey & Borrowed \\
\bottomrule
\end{tabular}
\end{center}

\subsection*{The non-parental share stays a single 2019 wave}

Pooling survey waves would cut the sampling error by 40 to 47 percent. I tested
it and rejected it, on the standard that every input matches the simulator's
2019 base year.

Pooling 2018, 2019 and 2020 centers exactly on 2019, which is the one pool that
passes the vintage test. It fails on content. NSCH 2020 was fielded through the
pandemic and child care use collapsed:

\begin{center}
\begin{tabular}{L{3.2cm}rrr}
\toprule
& \textbf{2018} & \textbf{2019} & \textbf{2020} \\
\midrule
Kentucky   & 57.3\% & 59.4\% & 45.8\% \\
New Jersey & 58.3\% & 60.2\% & 42.6\% \\
New York   & 64.7\% & 57.7\% & 62.1\% \\
Virginia   & 61.7\% & 52.0\% & 46.7\% \\
\bottomrule
\end{tabular}
\end{center}

That is a real shock, not sampling noise, and pooling it in would push every
state's level down by several points. Pooling 2018 and 2019 alone centers on
mid-2018 while every other input is pinned to 2019.

So the rate keeps its full sampling error. Both extra waves are held in the
repository so the rejection is reproducible, not so the pool can be used.

\subsection*{The borrowed fill rates are too high, measurably}

Two states have published their own figures and both point the same way.

Kentucky measured its centers at 57 percent full in 2019, against the national
90 percent that New York and Virginia borrow.

Virginia measured something related in 2022: what providers say they can
actually serve, against what they are licensed for. Centers came in at 73
percent of authorized capacity and homes at 91 percent, across 1,968 providers.
In that sample, 121,260 licensed slots against 89,374 children providers said
they could take.

Neither figure is used as a rate. Kentucky's applies to Kentucky. Virginia's is
a 2022 survey and measures capacity rather than enrollment, so it would
double-count against the fill rate. Both are evidence about direction and size,
and they agree: licensed capacity overstates real supply by roughly a quarter to
a third, and the national fill rate does not remove enough of it.

\section{What this method does not do}

\begin{enumerate}
\setlength{\itemsep}{0.4em}
\item \textbf{It cannot see unpaid relative care}, the largest informal
category. That number is a remainder, not a measurement, and no administrative
source anywhere would fix it. Only a restricted-access survey geocode file
would.

\item \textbf{The mix carries the register's date.} The level is 2019 by
construction, but how it divides across care types reflects whenever the state
last published. Composition genuinely drifts, and a 2026 register understates
family homes for 2019.

\item \textbf{Unpaid Center-Based is an upper bound} in every state, because
Head Start and state pre-K overlap by an amount nobody publishes, and a program
site may also hold a state licence.

\item \textbf{The scaling in step 5 forces an exact match} to the survey. That
is a choice rather than a finding, and it is worth revisiting: the alternative
is to let the registers say what they say and report the disagreement.

\item \textbf{The simulator still expects seven rows}, with center care split by
price. It rejects a file with six. Bridging the two is a mechanical step,
dividing the center row on a national share, and it is deliberately outside this
method so that the build does not appear to measure something it cannot.
\end{enumerate}

\end{document}
